\documentclass[12pt,a4paper]{article}

\usepackage{amsmath, amssymb, esvect, mathtools}
\newcommand{\R}{\mathbb{R}}
\newcommand{\N}{\mathbb{N}}
\newcommand{\e}{\mathrm{e}}
\newcommand{\vect}[1]{\vv{#1}}

\usepackage[utf8]{inputenc}
\usepackage{fontspec}
\usepackage{libertinus-otf}
\usepackage[french]{babel}
\usepackage{microtype}
\usepackage{icomma}

\setlength{\parindent}{0pt}
\usepackage[
    top=1.6cm,
    bottom=1.6cm,
    left=1.6cm,
    right=1.6cm,
    headheight=15pt,
    headsep=12pt,
    footskip=22pt
]{geometry}

\usepackage{fancyhdr}
\usepackage{lastpage}
\usepackage{tabularx}
\usepackage{tcolorbox}
\usepackage{enumitem}
\usepackage{eso-pic}
\usepackage{tikz}
\usepackage{xcolor}

\setlength{\abovedisplayskip}{6pt}
\setlength{\belowdisplayskip}{6pt}
\setlength{\abovedisplayshortskip}{3pt}
\setlength{\belowdisplayshortskip}{3pt}

\newcolumntype{Y}{>{\centering\arraybackslash}X}

\newcommand{\exercice}[3][5]{%
    \noindent{\Large\textbf{Exercice #2}} \hfill \textbf{#1 points}\\[0.20cm]%
    {\normalsize\textit{#3}}%
    \par\vspace{0.28cm}\hrule\vspace{0.55cm}%
}
\newcommand{\partie}[1]{\vspace{0.45cm}\noindent\textbf{\large Partie #1}\par\vspace{0.25cm}}

\setlist[enumerate,1]{label=\textbf{\arabic*.}, leftmargin=*, itemsep=0.35em, parsep=0pt, topsep=0.2em}
\setlist[enumerate,2]{label=\textbf{\alph*.}, leftmargin=*, itemsep=0.25em, parsep=0pt, topsep=0.15em}

\pagestyle{fancy}
\fancyhf{}
\fancyhead[L]{\small\textit{Bac Blanc -- Sujet 1}}
\fancyhead[C]{\small\textit{Session 2026}}
\fancyhead[R]{\small\textbf{Term\textsuperscript{spé}}}
\fancyfoot[C]{\thepage\ / \pageref{LastPage}}
\renewcommand{\headrulewidth}{0.4pt}
\renewcommand{\footrulewidth}{0.4pt}

\begin{document}

\thispagestyle{empty}

\AddToShipoutPictureFG*{%
    \begin{tikzpicture}[remember picture, overlay]
        \node[
            rotate=32,
            scale=1.1,
            text=black,
            opacity=0.06,
            font=\fontsize{72pt}{72pt}\bfseries\sffamily
        ] at (current page.center) {SP{\'E}CIALIT{\'E}};
    \end{tikzpicture}%
}

\noindent
\begin{minipage}[t]{0.52\textwidth}
    \textbf{\Large LYCÉE JEAN ROSTAND}\\[0.12cm]
    \textsc{Année scolaire 2025--2026}\\
    \textit{Classe de Terminale Générale}
\end{minipage}
\hfill
\begin{minipage}[t]{0.45\textwidth}
    \raggedleft
    \textbf{\large BACCALAURÉAT BLANC}\\[0.10cm]
    \textbf{Épreuve de spécialité}\\[0.10cm]
    \textit{Session 2026}\\[0.12cm]
    \small Date : Avril 2026
\end{minipage}

\vspace{0.30cm}
\hrule height 0.8pt
\vspace{0.45cm}

\vfill

\begin{center}
    {\LARGE \textbf{BACCALAURÉAT BLANC}}\\[0.20cm]
    {\large \textbf{SESSION 2026}}\\[0.15cm]
    {\textbf{\Large ÉPREUVE DE MATHÉMATIQUES}}\\[0.45cm]

    \begin{tcolorbox}[colback=white, colframe=black, arc=0mm, boxrule=1.5pt, width=0.96\textwidth, top=4mm, bottom=4mm]
        \centering
        {\fontsize{26pt}{30pt}\selectfont\textbf{MATHÉMATIQUES}}\\[0.30cm]
        {\Large \textsc{Épreuve d'Enseignement de Spécialité}}
    \end{tcolorbox}

    \vspace{0.50cm}
    {\large \textbf{Durée de l'épreuve : 4 heures}}\\[0.25cm]
    \textbf{L'usage de la calculatrice avec \underline{mode examen actif} est autorisé.}\\[0.10cm]
    \textit{L'usage de la calculatrice sans mémoire, « type collège », est autorisé.}
\end{center}

\vfill

\begin{center}
\renewcommand{\arraystretch}{1.5}
\begin{tabularx}{0.96\textwidth}{|>{\centering\arraybackslash\bfseries}m{2.6cm}|*{4}{Y|}>{\centering\arraybackslash\bfseries}m{2.2cm}|}
    \hline
    Exercice & \textbf{Exercice 1} & \textbf{Exercice 2} & \textbf{Exercice 3} & \textbf{Exercice 4} & \textbf{Total} \\ \hline
    Thème    & Probabilités & Analyse & Suites & Géométrie & --- \\ \hline
    Barème   & 4 points & 5 points & 6 points & 5 points & 20 points \\ \hline
\end{tabularx}
\end{center}

\vfill

\begin{center}
\begin{tcolorbox}[colback=gray!5, colframe=black, boxrule=1pt, arc=1mm, width=0.96\textwidth, top=3mm, bottom=3mm, left=4mm, right=4mm]
    \centering
    \textbf{\textsc{Consigne importante aux surveillants}}\\[0.15cm]
    \normalsize
    Nous remercions les surveillants de veiller à ce que chaque candidat inscrive lisiblement sur sa copie le \textbf{nom de son professeur de spécialité}.
\end{tcolorbox}
\end{center}

\vfill

\begin{tcolorbox}[colback=white, colframe=black, arc=0mm, boxrule=0.8pt, title=\textbf{\textsc{Consignes aux candidats}}, coltitle=black, colbacktitle=gray!15, top=3mm, bottom=3mm, left=4mm, right=4mm]
    \small
    \begin{itemize}[leftmargin=1.5em, itemsep=0.25em, label=---]
        \item Le sujet comporte \textbf{\pageref{LastPage} pages} numérotées de 1 à \pageref{LastPage}.
        \item Le candidat doit traiter les \textbf{4 exercices}.
        \item La qualité de la rédaction, la clarté et la précision des raisonnements seront prises en compte dans l'appréciation des copies.
        \item Tous les résultats doivent être justifiés, sauf mention contraire explicite.
        \item L'usage de tout document, téléphone ou appareil électronique connecté autre que la calculatrice autorisée est strictement interdit.
    \end{itemize}
\end{tcolorbox}

\newpage

\exercice[4]{1}{Thème : Probabilités et variables aléatoires}

Pour chacune des affirmations suivantes, indiquer si elle est vraie ou fausse. Chaque réponse devra être soigneusement justifiée ; une réponse non justifiée ne rapporte aucun point. Une absence de réponse ou une réponse inexacte n'enlève pas de point. Les questions de cet exercice sont indépendantes.

\vspace{0.35cm}

\begin{enumerate}
    \item Soient $A$ et $B$ deux événements de probabilités non nulles et indépendantes tels que $P(A) = 0{,}4$ et $P(A \cup B) = 0{,}7$.\\[0.1cm]
    \textbf{Affirmation 1 :} La probabilité de l'événement $B$ est égale à $0{,}5$.
    
    \vspace{0.55cm}
    \item Soit $X$ une variable aléatoire suivant la loi binomiale $\mathcal{B}(10; 0{,}3)$.\\[0.1cm]
    \textbf{Affirmation 2 :} La probabilité $P(X \ge 1)$ est supérieure ou égale à $0{,}97$.
    
    \vspace{0.55cm}
    \item Soit $Z$ une variable aléatoire suivant la loi normale centrée réduite $\mathcal{N}(0;1)$.\\[0.1cm]
    \textbf{Affirmation 3 :} $P(-1{,}96 \le Z \le 1{,}96) = 0{,}95$ à $10^{-2}$ près.
    
    \vspace{0.55cm}
    \item Soit $Y$ une variable aléatoire suivant une loi exponentielle de paramètre $\lambda > 0$ telle que $P(Y \ge 2) = \mathrm{e}^{-1}$.\\[0.1cm]
    \textbf{Affirmation 4 :} Le paramètre $\lambda$ est égal à $0{,}5$.
\end{enumerate}

\newpage

\exercice[5]{2}{Thème : Étude d'une fonction logarithmique et calcul intégral}

On considère la fonction $g$ définie sur l'intervalle $]0; +\infty[$ par :
\[
    g(x) = x\ln(x) - x\,.
\]
On note $\mathcal{C}_g$ la courbe représentative de $g$ dans un repère orthogonal.

\partie{A --- Étude de la fonction $g$}

\begin{enumerate}
    \item Déterminer la limite de la fonction $g$ en $0$.
    \item Déterminer la limite de la fonction $g$ en $+\infty$.
    \item Calculer $g'(x)$ pour tout $x \in ]0; +\infty[$ et dresser le tableau de variations de $g$.
    \item Démontrer que l'équation $g(x) = 0$ admet une unique solution $\alpha$ sur l'intervalle $[1; 3]$ et donner un encadrement d'amplitude $10^{-2}$ de $\alpha$.
\end{enumerate}

\partie{B --- Calcul d'une aire géométrique}

\begin{enumerate}
    \item Démontrer que la fonction $G$ définie sur $]0; +\infty[$ par $G(x) = \frac{1}{2}x^2\ln(x) - \frac{3}{4}x^2$ est une primitive de $g$ sur $]0; +\infty[$.
    \item Calculer l'aire exacte $\mathcal{A}$, en unités d'aire, du domaine plan délimité par l'axe des abscisses, la courbe $\mathcal{C}_g$ et les droites d'équations $x = 1$ et $x = \mathrm{e}$.
\end{enumerate}

\newpage

\exercice[6]{3}{Thème : Suites récurrentes et algorithmique Python}

On considère la suite $(u_n)$ définie par $u_0 = 2$ et, pour tout entier naturel $n$ :
\[
    u_{n+1} = \frac{3u_n + 2}{u_n + 4}\,.
\]

\partie{A --- Étude d'une fonction associée}

Soit $f$ la fonction définie sur $[0; +\infty[$ par $f(x) = \frac{3x+2}{x+4}$.

\begin{enumerate}
    \item Calculer $f'(x)$ et étudier les variations de $f$ sur $[0; +\infty[$.
    \item Démontrer que l'intervalle $[0; 2]$ est stable par $f$.
    \item 
    \begin{enumerate}
        \item Résoudre l'équation $f(x) = x$ sur $[0; 2]$.
        \item En déduire le point fixe $\ell$.
    \end{enumerate}
    \item Démontrer par récurrence que pour tout $n \in \mathbb{N}$, $0 \le u_n \le 2$.
\end{enumerate}

\partie{B --- Monotonie, limite et seuil algorithmique}

\begin{enumerate}
    \item Démontrer que pour tout $n \in \mathbb{N}$, $u_{n+1} - u_n = \frac{(1-u_n)(u_n+2)}{u_n+4}$.
    \item En déduire le sens de variation de la suite $(u_n)$.
    \item Justifier que la suite $(u_n)$ converge et déterminer sa limite.
    \item On souhaite déterminer à l'aide d'un script Python le rang à partir duquel $u_n$ est proche de $1$ à $\varepsilon$ près.
\end{enumerate}

\newpage

\exercice[5]{4}{Thème : Géométrie dans l'espace et produit scalaire}

Dans l'espace muni d'un repère orthonormé $(O; \vec{i}, \vec{j}, \vec{k})$, on considère les points :
\[
    A(1;0;2), \quad B(0;2;1), \quad C(2;1;0) \quad \text{et} \quad D(3;3;3)\,.
\]

\partie{A --- Plans et tétraèdre}

\begin{enumerate}
    \item Démontrer que les points $A, B$ et $C$ définissent un plan $\mathcal{P}$.
    \item Démontrer qu'une équation cartésienne du plan $\mathcal{P}$ est : $x + y + z - 3 = 0$.
    \item Vérifier que le point $D$ n'appartient pas au plan $\mathcal{P}$.
    \item Calculer le volume du tétraèdre $ABCD$.
\end{enumerate}

\partie{B --- Orthogonalité, projeté et sphère}

\begin{enumerate}
    \item Déterminer une représentation paramétrique de la droite $\Delta$ passant par $D$ et orthogonale au plan $\mathcal{P}$.
    \item Déterminer les coordonnées du projeté orthogonal $H$ du point $D$ sur le plan $\mathcal{P}$.
    \item Calculer la distance $DH$.
    \item Soit $\mathcal{S}$ la sphère de centre $D$ et de rayon $R = 3$. Étudier l'intersection de la sphère $\mathcal{S}$ avec le plan $\mathcal{P}$.
\end{enumerate}

\end{document}